\begin{figure}[htbp]
			\centering
			\resizebox{\textwidth}{!}{%
				\begin{tikzpicture}[
					scale=0.95,
					>=latex,
					trave/.style={line width=1.2pt},
					carico/.style={black, line width=0.4pt},
					forza/.style={black, line width=1pt},
					diagT/.style={black, line width=0.8pt},
					diagM/.style={black, line width=0.8pt},
					hatch/.style={gray!50, line width=0.3pt},
					quota/.style={black, line width=0.4pt},
					fibre/.style={gray!70, line width=0.3pt},
					]
					
					% ── Parametri dei vincoli (convenzione C06/C07) ──
					\def\tw{0.15}    \def\thR{0.28}   \def\rn{0.06}
					\def\rw{0.05}    \def\bw{0.09}    \def\gw{0.22}
					\def\hh{0.10}
					
					% ── Parametri dell'incastro (muro verticale + campitura terra) ──
					\def\twall{0.18}    % larghezza della campitura del muro a sinistra di A
					\def\hwall{0.36}    % semi-altezza del muro
					
					% ── Geometria del problema ──
					\def\Lell{4}                            % lunghezza trave ℓ
					\pgfmathsetmacro{\xA}{0}
					\pgfmathsetmacro{\xMid}{0.5*\Lell}      % mezzeria (per posizionare l'etichetta della quota)
					\def\yqi{-0.80}                         % y delle quote sotto la trave
					\def\ylab{-3.30}                        % y delle etichette di pannello
					
					% ── Parametri del carico distribuito ──
					\def\pgap{0.05}                         % gap fra trave e base del riquadro
					\def\php{0.40}                          % altezza del riquadro
					\pgfmathsetmacro{\pBase}{\pgap}         % base del riquadro: y = 0.05
					\pgfmathsetmacro{\pTop}{\pgap+\php}     % cima del riquadro: y = 0.45
					
					% =========================================================
					% PANNELLO (a) — Schema statico
					% =========================================================
					\begin{scope}[xshift=0cm]
						
						% ── Trave ──
						\draw[trave] (0,0) -- (\Lell,0);
						
						% ── Fibra inferiore (tratteggio sotto la trave) ──
						\draw[fibre, dashed] (0, -0.1) -- (\Lell, -0.1);
						
						% ── Incastro in A: muro verticale + campitura "terra" a 45° ──
						% Linea verticale del muro
						\draw[thick] (0,-\hwall) -- (0,\hwall);
						% Tratti obliqui (campitura), equispaziati lungo y
						\foreach \y in {-0.30,-0.18,-0.06,0.06,0.18,0.30} {
							\draw[thin] (0,\y) -- ++(-\twall,-\twall);
						}
						
						% ── Etichette A, B ──
						\node[below right=-2pt, font=\small] at (0,0)      {$A$};
						\node[below right=-2pt, font=\small] at (\Lell,0)  {$B$};
						
						% ── Carico distribuito p̄ sopra la trave ──
						\draw[carico] (0,\pBase) -- (\Lell,\pBase);
						\draw[carico] (0,\pTop)  -- (\Lell,\pTop);
						\foreach \x in {0, 0.50, 1.00, ..., 4.00} {
							\draw[->, carico, >={Latex[length=2pt, width=1.5pt]}]
							(\x, \pTop) -- (\x, \pBase);
						}
						\node[above, font=\small] at (\Lell/2, \pTop) {$\bar{p}$};
						
						% ── Quota ℓ sotto la trave ──
						\draw[<->, quota] (0,\yqi) -- (\Lell,\yqi);
						\draw[quota] (0,\yqi+0.07)     -- (0,\yqi-0.07);
						\draw[quota] (\Lell,\yqi+0.07) -- (\Lell,\yqi-0.07);
						\node[above, font=\small] at (\xMid, \yqi-0.04) {$\ell$};
						
						% Etichetta pannello (a)
						\node[font=\small] at (0.5*\Lell, 0.4*\ylab) {(a)};
						
					\end{scope}
					
					% =========================================================
					% PANNELLO (b) — Diagramma di T(z)
					% =========================================================
					\begin{scope}[xshift=5.5cm]
						
						% ── Trave (clone) ──
						\draw[trave] (0,0) -- (\Lell,0);
						
						% ── Incastro in A (clone) ──
						\draw[thick] (0,-\hwall) -- (0,\hwall);
						\foreach \y in {-0.30,-0.18,-0.06,0.06,0.18,0.30} {
							\draw[thin] (0,\y) -- ++(-\twall,-\twall);
						}
						
						% ── Etichette A, B ──
						\node[below right=-2pt, font=\small] at (0,0)      {$A$};
						\node[below right=-2pt, font=\small] at (\Lell,0)  {$B$};
						
						% ── Diagramma T(z) ──
						\def\sT{0.25}
						\pgfmathsetmacro{\TA}{\Lell*\sT}        % T(A) = +p̄ℓ → +0.80
						
						% Hatching verticale dall'asse al diagramma
						\foreach \x in {0.20, 0.40, ..., 3.80} {
							\pgfmathsetmacro{\h}{\sT*(\Lell-\x)}
							\draw[hatch] (\x, 0) -- (\x, \h);
						}
						
						% Bordo del diagramma — triangolo: verticale in A, retta fino a B
						\draw[diagT]
						(0, 0) -- (0, \TA)              % verticale di chiusura in A
						-- (\Lell, 0);                  % lineare da +T_A in A a 0 in B
						
						% Segno cerchiato nel triangolo positivo
						\node[font=\small] at ({\Lell/3}, {\TA/3}) {$\boldsymbol{\oplus}$};
						
						% Valore notevole — etichetta sopra il vertice in A
						\node[left, font=\small] at (0, \TA) {$\bar{p}\ell$};
						
						% Etichetta pannello (b)
						\node[font=\small] at (0.5*\Lell, 0.4*\ylab) {(b)};
						
					\end{scope}
					
					% =========================================================
					% PANNELLO (c) — Diagramma di M(z)
					% =========================================================
					\begin{scope}[xshift=11.0cm]
						
						% ── Trave (clone) ──
						\draw[trave] (0,0) -- (\Lell,0);
						
						% ── Incastro in A (clone) ──
						\draw[thick] (0,-\hwall) -- (0,\hwall);
						\foreach \y in {-0.30,-0.18,-0.06,0.06,0.18,0.30} {
							\draw[thin] (0,\y) -- ++(-\twall,-\twall);
						}
						
						% ── Etichette A, B ──
						\node[above right=-2pt, font=\small] at (0,0)      {$A$};
						\node[above right=-2pt, font=\small] at (\Lell,0)  {$B$};
						
						% ── Diagramma M(z) ──
						\def\sM{0.15}
						
						% Hatching verticale: dall'asse al diagramma (sopra l'asse)
						\foreach \x in {0.20, 0.40, ..., 3.80} {
							\pgfmathsetmacro{\h}{\sM*(\Lell-\x)*(\Lell-\x)/2}
							\draw[hatch] (\x, 0) -- (\x, \h);
						}
						
						% Bordo del diagramma — parabola campionata a passo 0.20
						\pgfmathsetmacro{\MA}{\sM*\Lell*\Lell/2}    % |M(A)| in figura = sM·ℓ²/2 = 0.80
						\draw[diagM]
						(0, 0) -- (0, \MA)
						\foreach \x in {0.20, 0.40, ..., 3.80} {
							-- ({\x}, {\sM*(\Lell-\x)*(\Lell-\x)/2})
						}
						-- (\Lell, 0);
						
						% Valore notevole — etichetta sopra il vertice in A
						\node[left, font=\small] at (0, \MA) {$\tfrac{\bar{p}\ell^2}{2}$};
						
						% Etichetta pannello (c)
						\node[font=\small] at (0.5*\Lell, 0.4*\ylab) {(c)};
						
					\end{scope}
					
				\end{tikzpicture}
			}
			\caption{Mensola con carico uniformemente distribuito
				$p(z) = \bar{p}$ (verso il basso), incastrata in $A$ e
				libera in $B$:
				(a)~schema strutturale;
				(b)~diagramma del taglio $T(z)$, lineare decrescente, con
				valore $\bar{p}\ell$ all'incastro e annullamento
				all'estremità libera;
				(c)~diagramma del momento flettente $M(z)$, parabolico
				con concavità verso l'alto, sempre negativo, con valore
				minimo $-\bar{p}\ell^2/2$ all'incastro e annullamento
				all'estremità libera, dove il diagramma arriva con
				tangente orizzontale.}
			\label{fig:stat_trasv_iso_03}
		\end{figure}
